
\newpage
\section{Теоретико-методологические основы построения агентных систем для анализа научных данных}

\subsection{Эволюция парадигм интеллектуальных систем от экспертных систем к мультиагентным архитектурам}

\subsubsection{Исторический анализ развития экспертных систем в научных исследованиях}

Эволюция искусственного интеллекта (ИИ) представляет собой сложный путь от попыток формализации логических рассуждений до создания систем, способных к автономным действиям. Двумя ключевыми вехами на этом пути стали экспертные системы, доминировавшие в 1970–1980-х годах, и современный агентно-ориентированный подход (Agentic AI). Анализ их развития позволяет не только проследить смену технологических парадигм, но и выявить фундаментальные принципы, которые остаются неизменными при создании интеллектуальных систем, способных решать сложные, слабо формализуемые задачи.

\subsubsection{Становление агентно-ориентированного подхода в искусственном интеллекте}

В данном анализе рассматривается исторический путь от первых экспертных систем, основанных на статичных правилах, к агентным системам, демонстрирующим целенаправленное рассуждение и автономное выполнение действий.

\subsubsection{Таксономия интеллектуальных агентов по функциональному назначению и архитектурным принципам}

Текст раздела...

\subsubsection{Современные программные фреймворки для конструирования мультиагентных систем}

Текст раздела...

\subsection{Большие языковые модели как фундаментальный компонент интеллектуальных систем}

\subsubsection{Архитектурные принципы и эволюция больших языковых моделей}

Текст раздела...

\subsubsection{Классификация и сравнительный анализ современных LLM для решения научных задач}

Текст раздела...

\subsubsection{Проприетарные модели: архитектура и особенности применения (OpenAI GPT, Anthropic Claude, Google Gemini, Cohere)}

Текст раздела...

\subsubsection{Открытые и исследовательские модели: особенности и преимущества (BLOOM, LLaMA, Mistral, Qwen)}

Текст раздела...

\subsubsection{Модели для локального развертывания: оптимизация и инференс (Local LLM, Custom Model)}

Текст раздела...

\subsubsection{Методология дообучения языковых моделей на предметных данных}

Текст раздела...

\subsubsection{Сравнительный анализ эффективности различных LLM в научных исследованиях}

Текст раздела...

\subsection{Теоретические основы извлечения информации из гетерогенных источников}

\subsubsection{Методы семантического поиска и векторного представления текстов}

Текст раздела...

\subsubsection{Архитектура RAG: принципы и компоненты}

Текст раздела...

\subsubsection{Графовые модели представления знаний в биоинформатике}

Текст раздела...

\subsubsection{Гибридные стратегии поиска информации в реляционных базах данных}

Текст раздела...

\subsection{Технологии мультимодального человеко-машинного взаимодействия}

\subsubsection{Современные методы автоматического распознавания речи}

Текст раздела...

\subsubsection{Технологии синтеза речи и их применение в научных системах}

Текст раздела...

\subsubsection{Интеграция текстовых, голосовых и графических интерфейсов}

Текст раздела...

\subsubsection{Методы обработки голосовых команд в специализированных областях}

Текст раздела...

\subsection{Мультимодальные диалоговые системы: архитектура и принципы}

\subsubsection{Концепция мультимодального диалога}

Текст раздела...

\subsubsection{Методы синхронизации различных модальностей ввода/вывода}

Текст раздела...

\subsubsection{Управление контекстом в мультимодальных системах}

Текст раздела...

\subsubsection{Оценка качества пользовательского опыта в мультимодальных интерфейсах}

Текст раздела...

\subsection{Лингвистическое обеспечение многоязычных интеллектуальных систем}

\subsubsection{Методы автоматической идентификации языка текстов}

Текст раздела...

\subsubsection{Технологии машинного перевода научной терминологии}

Текст раздела...

\subsubsection{Многоязычные языковые модели и их применение}

Текст раздела...

\subsubsection{Геолокационные сервисы и адаптация контента}

Текст раздела...

\subsubsection{Многоязычная поддержка в голосовых интерфейсах}

Текст раздела...

\subsection{Методы компьютерного анализа химических данных}

\subsubsection{Формальные языки описания химических структур}

Текст раздела...

\subsubsection{Методы расчета молекулярных дескрипторов}

Текст раздела...

\subsubsection{Компьютерное зрение для распознавания химических структур из изображений}

Текст раздела...

\subsubsection{Интеграция химических инструментов с LLM}

Текст раздела...

\subsubsection{Голосовой ввод химических формул и названий}

Текст раздела...

\subsection{Хронобиология как предметная область интеллектуального анализа}

\subsubsection{Фундаментальные концепции хронобиологии и хронофармакологии}

Текст раздела...

\subsubsection{Уникальность базы данных хронобиотиков как единственного в мире структурированного источника знаний в предметной области}

Текст раздела...

\subsubsection{Анализ структуры и содержания базы данных хронобиотических веществ}

Текст раздела...

\subsubsection{Специфика анализа биологических ритмов и временных зависимостей}

Текст раздела...

\subsubsection{Терминологические особенности предметной области для голосового ввода}

Текст раздела...

\subsubsection{Научная и практическая значимость автоматизации доступа к уникальным данным}

Текст раздела...

\subsection{Архитектурные паттерны проектирования специализированных агентов}

\subsubsection{Принципы функциональной декомпозиции в агентных системах}

Текст раздела...

\subsubsection{Методы координации и коммуникации между агентами}

Текст раздела...

\subsubsection{Шаблоны проектирования реактивных и когнитивных агентов}

Текст раздела...

\subsubsection{Организация мультимодального взаимодействия на уровне агентов}

Текст раздела...

\subsubsection{Интеграция инструментов расширения функциональности}

Текст раздела...

\subsection{Методология поиска и анализа научной литературы}

\subsubsection{Современные методы автоматизированного поиска научных публикаций}

Текст раздела...

\subsubsection{Интеграция с академическими базами данных}

Текст раздела...

\subsubsection{Методы извлечения ключевой информации из научных статей}

Текст раздела...

\subsubsection{Обнаружение новых публикаций и отслеживание научных трендов}

Текст раздела...

\subsection{Методология оценки эффективности интеллектуальных систем}

\subsubsection{Количественные метрики качества извлечения информации}

Текст раздела...

\subsubsection{Методы валидации результатов на ограниченных выборках}

Текст раздела...

\subsubsection{Критерии качества работы специализированных агентов}

Текст раздела...

\subsubsection{Оценка эффективности мультимодального взаимодействия}

Текст раздела...

\subsection{Анализ существующих решений и постановка задачи исследования}

\subsubsection{Критический анализ существующих систем анализа научных данных}

Текст раздела...

\subsubsection{Идентификация ограничений существующих подходов}

Текст раздела...

\subsubsection{Обоснование необходимости создания интеллектуальной системы для единственной в мире базы хронобиотиков}

Текст раздела...

\subsubsection{Формальная постановка задачи разработки мультимодальной агентной системы}

Текст раздела...

\subsection{Выводы и результаты по главе}

Текст раздела...